% =============================================================================
% Section 9: Conclusion
% =============================================================================

We have presented a two-parameter model that determines the sampling probability for network centrality approximation from network reach alone. On our synthetic suite, a power-law fit yields an exponent that is consistent with square-root scaling (we do not reject $\alpha = 0.5$), and model comparison across five candidate functional forms supports the robustness of the proposed specification. Validation on the Greater London (\glaNnodes{} nodes) and Madrid street networks achieves Spearman $\rho > \glaMinRho$ at all tested distances (5--20km), with speedups of \glaFiveKmSpeedup{}--\glaTwentyKmSpeedup{}$\times$, demonstrating that the model extrapolates well beyond the synthetic networks on which it was trained.

By adapting the Eppstein--Wang source-sampling bound to the localised setting, we find that the resulting Hoeffding-style bound holds for the rank model's recommended sample sizes in over 93\% of synthetic configurations, with most violations confined to a low-reach regime where sampling provides limited benefit and concentration guarantees are weak. At high reach, the two perspectives become numerically consistent: under the localised bound analysis, rank-calibrated sample sizes correspond to an implied normalised additive error on the order of $\varepsilon \approx 0.1$. The rank model and the localised error bound thus address complementary needs: rank preservation for comparative urban analysis, and additive accuracy for applications requiring absolute centrality estimates.

The practical value lies in providing analysts with a calibrated way to trade computation time for approximation error, with evidence that sampled results preserve ranking accuracy and, at large reach, correspond to tighter implied additive-error levels under a conservative concentration analysis. The largest benefits accrue at long distances on large networks, precisely where computational cost is highest. More broadly, our results suggest that the effective sample size--accuracy relationship is not purely an artifact of the synthetic configurations studied here, though further validation on non-European cities would strengthen this claim.

The model is implemented in the open-source \cityseer{} Python library.\footnote{\texttt{https://github.com/benchmark-urbanism/cityseer-api}} Fitting scripts, validation data, and model comparison analyses are available in the repository's \texttt{analysis/sampling/} directory.
